\section{Introducción}

El confinamiento espacial de la materia a la nanoescala rompe la simetría
traslacional del estado sólido masivo, originando propiedades optofísicas
emergentes \cite{lasmiSilverNanoparticlesAgNPs2025}. En este régimen, las
nanopartículas de plata (AgNPs) constituyen un sistema físico de interés
excepcional. Su estudio trasciende la ciencia de materiales básica,
fundamentando el desarrollo de tecnologías que abarcan desde la biomedicina
avanzada hasta la optoelectrónica \cite{hosnyComprehensiveReviewSilver2025,
satiSilverNanoparticlesAgNPs2025}.

La electrodinámica de las AgNPs está dominada por la resonancia de plasmón de
superficie localizado (LSPR). Este fenómeno emerge del acoplamiento fuerte entre
un campo electromagnético incidente y las oscilaciones coherentes del plasma de
electrones de conducción \cite{lasmiSilverNanoparticlesAgNPs2025}. La
condición de resonancia confiere a las AgNPs secciones eficaces de absorción
y dispersión extremadamente altas, además de una sensibilidad rigurosa a las
fluctuaciones de la permitividad dieléctrica del entorno. Esta dependencia
sustenta su implementación en biosensores de resolución molecular y aplicaciones
teranósticas \cite{chaionSilverNanoparticleBased2025,
satiSilverNanoparticlesAgNPs2025}.

En el régimen de campos electromagnéticos intensos, el confinamiento local
amplifica la susceptibilidad óptica no lineal del sistema. Evaluadas
experimentalmente mediante mediciones de escaneo Z, las AgNPs exhiben
dinámicas transitorias complejas como la absorción saturable y la absorción
inversa saturable \cite{qayyumNonlinearOpticalProperties2025}. El control
determinista de estas propiedades, requisito ineludible para el diseño de
limitadores ópticos y dispositivos de procesamiento fotónico, es una función
estricta de la topología geométrica y la estabilidad termodinámica del sistema.

Clásicamente, la reducción de iones argénticos ha dependido de vías químicas
que involucran precursores de alta toxicidad. La transición hacia metodologías
biogénicas representa una optimización física y ambiental. La utilización de
extractos biológicos provee un medio de reacción donde los fitocompuestos actúan
simultáneamente como agentes reductores y pasivantes
\cite{thomasPlantbasedSynthesisCharacterization2024,
satiSilverNanoparticlesAgNPs2025}. Esta biocorona no solo minimiza la
citotoxicidad para viabilizar la experimentación \textit{in vivo}, sino que
redefine las condiciones de frontera dieléctricas de la partícula, alterando de
forma directa la dinámica plasmónica
\cite{hosnyComprehensiveReviewSilver2025}.

A pesar del volumen de literatura disponible, resulta imperativo establecer un
marco formal que integre las metodologías de síntesis sostenibles con la
electrodinámica no lineal fundamental. El objetivo de esta revisión sistemática
es estructurar el conocimiento reciente (2020--2026) sobre la física de las
AgNPs. Se prioriza el análisis de la transición hacia vías de síntesis verdes,
su impacto sobre la sintonización de la LSPR, y la consecuente modulación de la
respuesta óptica no lineal, identificando las barreras teóricas y experimentales
para la ingeniería de dispositivos optoelectrónicos escalables.
